\section{Discusión}

\pendiente{Esta sección se termina con los datos de la Etapa~H. Lo que sigue es el esqueleto
con las diferencias ya medidas entre el modelo y la planta. Ordenarlas por peso cuando se tengan
los resultados finales y no corregirlas artificialmente.}

La diferencia principal entre el modelo y el prototipo está en la transmisión entre la orden de
inclinación y la bola. El modelo de diseño suponía $c = 5/7$ y un retardo de
\SI{75}{\milli\second}. La identificación dio una ganancia efectiva cercana a \num{0.2} para
movimientos pequeños y un retardo de \SIrange{0.27}{0.41}{\second}, con una zona muerta de
alrededor de \SI{1}{\degree}. Como el resultado fue el mismo con dos bolas muy distintas, el
origen está en el mecanismo: holgura de las rótulas y varillas, medida hasta \SI{2.4}{\degree}
en el eje Y, y respuesta del servo a cambios pequeños. El eje Y, cuyo servo es más duro, tiene
más holgura y más retardo, y es donde aparecieron casi todas las zonas en las que la bola se
quedaba trabada. La asimetría entre los ejes viene, entonces, del actuador.

Estas diferencias explican el comportamiento observado. Una ganancia de planta cuatro veces
menor queda cerca del margen de ganancia inferior de \SI{-12.6}{\deci\bel} del lazo
condicionalmente estable (Tabla~\ref{tab:controlador}), y un retardo cuatro veces mayor deja muy
poco margen de fase. Con un retardo total estimado de \SI{0.44}{\second} (cámara, datos de la
ESP32 y transmisión), el cálculo de márgenes del controlador original da un margen de fase de
unos \SI{12}{\degree} para $c_{\mathrm{ef}} = 0{,}2$, y un rediseño preliminar con ese modelo
($\Kp = \SI{0.00727}{\degree\per\milli\metre}$, $\Kd = \SI{0.0184}{\degree\per(\milli\metre\per\second)}$,
$\Ti = \SI{5}{\second}$) sube el margen a \SI{43}{\degree}. Con ese retardo, el requisito R1 de
\SI{3}{\second} no se puede cumplir con un controlador robusto: o se acepta un requisito menos
exigente o se reduce el retardo, con menos holgura mecánica y un \texttt{loop()} más rápido en la
ESP32.

\pendiente{Los márgenes del rediseño se calcularon fuera de MATLAB (Python, sesión del
2026-11-13). Antes de citarlos, repetir C1 y D1 en MATLAB con el modelo nominal identificado y
regenerar el código (E a G). Comparar en el prototipo el original y el rediseño con el mismo
ensayo: en las pruebas del 2026-11-13 el rediseño terminó más cerca en los escalones (error final
de \SIrange{1}{47}{\milli\metre}, contra \SIrange{9}{80}{\milli\metre}), pero siguió peor el
círculo (error mediano de \SIrange{45}{58}{\milli\metre}, contra \SI{30}{\milli\metre}).}

Otras diferencias identificadas: la latencia medida de la visión fue menor que la supuesta, el
ruido de posición fue unas diez veces menor que el \SI{1}{\milli\metre} del diseño (aunque, en
simulación, bajar el parámetro del Kalman a ese valor hacía oscilar el lazo, por lo que se
mantuvo en \SI{1}{\milli\metre}), la relación servo--plato fue asimétrica y más empinada que la calculada, y
el término centrífugo limita la validez del modelo lineal cuando el plato se mueve rápido.

\pendiente{Diferencias de la implementación física (cuantización, efecto del \texttt{loop()} de
\SI{40}{\milli\second}, la zona de tolerancia de \SI{5}{\milli\metre} y su efecto sobre R3,
perturbaciones por choques con el borde) cuando se tengan los datos de H.}
